\begin{proof}
\textbf{Degenerate and edge cases.}\\[4pt]
\emph{Case \(n=1\).}  
The hypothesis gives \(\lambda_{1}=1\); hence the statement reads  
\(f(x_{1})\le f(x_{1})\), which holds with equality.

\emph{Zero weights.}  
If \(\lambda_{k}=0\) for some index \(k\), then the term \(\lambda_{k}x_{k}\) (resp. \(\lambda_{k}f(x_{k})\)) disappears from the left‑hand (resp. right‑hand) side.  Thus the inequality for \(\{x_{i},\lambda_{i}\}_{i=1}^{n}\) is equivalent to the inequality for the reduced family obtained by omitting \(k\).  Hence we may assume \(\lambda_{i}>0\) for all \(i\).

\medskip
\textbf{The case \(n=2\).}\\[4pt]
Let \(\lambda_{1},\lambda_{2}\ge0\) with \(\lambda_{1}+\lambda_{2}=1\).  By convexity,
\[
f(t x+(1-t)y)\le t\,f(x)+(1-t)f(y)\qquad (t\in[0,1],\;x,y\in\mathbb R).
\]
Choosing \(t=\lambda_{1}\), \(x=x_{1}\), \(y=x_{2}\) yields
\[
f(\lambda_{1}x_{1}+\lambda_{2}x_{2})\le
\lambda_{1}f(x_{1})+\lambda_{2}f(x_{2}),
\]
which is Jensen’s inequality for two points.

\medskip
\textbf{Inductive hypothesis.}\\[4pt]
Assume that for some integer \(k\ge2\) and for every choice of points
\(y_{1},\dots ,y_{k}\) and weights \(\mu_{1},\dots ,\mu_{k}\) with \(\mu_{i}\ge0\) and
\(\sum_{i=1}^{k}\mu_{i}=1\) we have
\begin{equation}
f\!\Bigl(\sum_{i=1}^{k}\mu_{i}y_{i}\Bigr)\le
\sum_{i=1}^{k}\mu_{i}f(y_{i}). \tag{IH}
\end{equation}

\medskip
\textbf{Inductive step \(k\to k+1\).}\\[4pt]
Let \(\lambda_{1},\dots ,\lambda_{k+1}\ge0\) satisfy \(\sum_{i=1}^{k+1}\lambda_{i}=1\).
Set
\[
\alpha:=\lambda_{k+1},\qquad
\beta:=1-\alpha=\sum_{i=1}^{k}\lambda_{i}.
\]
If \(\beta=0\) then \(\alpha=1\) and the desired inequality reduces to
\(f(x_{k+1})\le f(x_{k+1})\).  Hence we may assume \(\beta>0\).

Define the intermediate point
\[
y:=\frac{1}{\beta}\sum_{i=1}^{k}\lambda_{i}x_{i}.
\]
The numbers \(\frac{\lambda_{i}}{\beta}\;(i=1,\dots ,k)\) are non‑negative and sum to \(1\), so \(y\) is a convex combination of the first \(k\) points.

\smallskip
\emph{4.1 Applying two‑point convexity.}  
Since \(\alpha\in[0,1]\) and \(\beta=1-\alpha\), convexity gives
\begin{equation}
f\bigl(\beta y+\alpha x_{k+1}\bigr)\le
\beta\,f(y)+\alpha\,f(x_{k+1}). \tag{1}
\end{equation}
Moreover,
\begin{align*}
\beta y+\alpha x_{k+1}
   &=\beta\Bigl(\frac{1}{\beta}\sum_{i=1}^{k}\lambda_{i}x_{i}\Bigr)
     +\lambda_{k+1}x_{k+1} \\
   &=\sum_{i=1}^{k}\lambda_{i}x_{i}+\lambda_{k+1}x_{k+1}
   =\sum_{i=1}^{k+1}\lambda_{i}x_{i}. \tag{2}
\end{align*}
Thus the left–hand side of \((1)\) is precisely the term we need to bound.

\smallskip
\emph{4.2 Estimating \(\beta f(y)\) via the induction hypothesis.}  
Applying \((\text{IH})\) to the points \(x_{1},\dots ,x_{k}\) with weights
\(\frac{\lambda_{i}}{\beta}\) yields
\begin{equation}
f(y)=f\!\Bigl(\sum_{i=1}^{k}\frac{\lambda_{i}}{\beta}x_{i}\Bigr)
     \le\sum_{i=1}^{k}\frac{\lambda_{i}}{\beta}\,f(x_{i}). \tag{3}
\end{equation}
Multiplying \((3)\) by \(\beta>0\) gives
\begin{equation}
\beta f(y)\le\sum_{i=1}^{k}\lambda_{i}f(x_{i}). \tag{4}
\end{equation}

\smallskip
\emph{4.3 Combining the estimates.}  
From \((1)\) and \((2)\) we have
\[
f\!\Bigl(\sum_{i=1}^{k+1}\lambda_{i}x_{i}\Bigr)
   \le \beta f(y)+\alpha f(x_{k+1}).
\]
Using \((4)\) and \(\alpha=\lambda_{k+1}\) we obtain
\[
f\!\Bigl(\sum_{i=1}^{k+1}\lambda_{i}x_{i}\Bigr)
   \le \sum_{i=1}^{k}\lambda_{i}f(x_{i})+\lambda_{k+1}f(x_{k+1})
   =\sum_{i=1}^{k+1}\lambda_{i}f(x_{i}),
\]
which completes the inductive step.

\medskip
\textbf{Conclusion of the induction.}  
The inequality holds for \(n=2\) (Section 2) and the step \(k\to k+1\) has been proved; therefore, by induction, Jensen’s inequality holds for every integer \(n\ge2\).  Together with the trivial case \(n=1\) (Section 1) we obtain the result for all \(n\in\mathbb N\).

\medskip
\textbf{Equality conditions.}\\[4pt]
\begin{enumerate}
\item If one weight equals \(1\) and the others are zero, both sides reduce to \(f(x_{j})\) for that index \(j\).
\item If \(f\) is affine on the convex hull of the points \(\{x_{i}\mid\lambda_{i}>0\}\), i.e. there exist \(a,b\in\mathbb R\) such that \(f(t)=at+b\) on that hull, then equality holds.
\item In all other situations at least one of the intermediate convexity inequalities used in the inductive step is strict, so the overall inequality is strict.
\end{enumerate}
\end{proof}